% =====================================================================
\section{Exercise 4 -- Option-based portfolios under alternative copulas}
\label{sec:ex4}
% =====================================================================

This exercise reproduces the Monte-Carlo experiment of the
\emph{Option-based portfolios} slides (pages 16--19) using a Python port of
\file{ccK\_simuRet2.m}. The full implementation is in \file{exercise4/}.

% ---------------------------------------------------------------
\subsection{Setup}
% ---------------------------------------------------------------

Two assets with log-normal terminal prices
\[
S_{iT}=S_{i0}\exp\!\Bigl(\bigl(\mu-\tfrac12\sigma^2\bigr)T+\sigma\sqrt{T}\,Z_i\Bigr),
\qquad i=1,2,
\]
where $S_{10}=S_{20}=100$, $r=0.03$, $\mu=0.08$, $\sigma=0.2$, $T=1$, and
$(Z_1,Z_2)$ has $\mathcal{N}(0,1)$ marginals coupled by one of four copulas
(scenarios 4.1--4.4 below). $N=10^5$ paths and a strike grid
$K\in\{85,86,\ldots,115\}$.

\paragraph{Black--Scholes prices (port of \file{BS.m}).}
\[
d_1=\frac{\ln(S_0/K)+\bigl(r+\tfrac12\sigma^2\bigr)T}{\sigma\sqrt T},\quad
d_2=d_1-\sigma\sqrt T,
\]
\[
C_0(K,T)=S_0\Phi(d_1)-Ke^{-rT}\Phi(d_2),\qquad
P_0(K,T)=Ke^{-rT}\Phi(-d_2)-S_0\Phi(-d_1).
\]

\paragraph{Portfolio returns.}
We consider four standard strategies (slides eqs.~4--5):
\begin{align}
R_{\text{cc}}(K,T) &=\frac{S_{1T}-\max(S_{1T}-K,0)}{S_{10}-C_0(K,T)}-1
&&\text{(covered call, single asset)}\\
R_{\text{cc2}}(K,T)&=\frac{S_{2T}-\max(S_{1T}-K,0)}{S_{20}-C_0(K,T)}-1
&&\text{(covered call, hedged with asset 1)}\\
R_{\text{pp}}(K,T) &=\frac{S_{1T}+\max(K-S_{1T},0)}{S_{10}+P_0(K,T)}-1
&&\text{(protective put, single asset)}\\
R_{\text{pp2}}(K,T)&=\frac{S_{2T}+\max(K-S_{1T},0)}{S_{20}+P_0(K,T)}-1
&&\text{(protective put, hedged with asset 1)}
\end{align}
\code{cc} and \code{pp} use only $S_{1T}$ and are therefore
\emph{copula-independent}; \code{cc2} and \code{pp2} mix $S_{1T}$ and $S_{2T}$
and so depend on the joint distribution. This is the central distinction the
slides emphasise.

\paragraph{Generation of $(Z_1,Z_2)$.}
\begin{itemize}\setlength\itemsep{0pt}
\item Gaussian copula: direct Cholesky $Z_1=\varepsilon_1$, $Z_2=\rho\varepsilon_1+\sqrt{1-\rho^2}\,\varepsilon_2$.
\item Student-$t$ copula: simulate $(U_1,U_2)\sim C^t_{\rho,\nu}$ via \eqref{eq:t-sim}
      then $Z_i=\Phi^{-1}(U_i)$.
\item Clayton copula: simulate $(U_1,U_2)\sim C^{Cla}_\theta$ via \eqref{eq:clayton-sim}
      then $Z_i=\Phi^{-1}(U_i)$.
\end{itemize}

% ---------------------------------------------------------------
\subsection{Results across scenarios}
% ---------------------------------------------------------------

Table~\ref{tab:4-atm} aggregates the at-the-money ($K=100$) statistics of each
strategy under the four copula scenarios. The full strike-curve plots are in
Figures~\ref{fig:4-1}--\ref{fig:4-4}.

\begin{table}[H]\centering\footnotesize
\caption{Strategy statistics at $K=100$, $N=10^5$ Monte-Carlo paths.}\label{tab:4-atm}
\begin{tabular}{llrrrr}
\toprule
Scenario & Strategy & Mean & Std & Skew & Kurt\\\midrule
\multirow{4}{*}{4.1 Gaussian $\rho=0.9$}
 & cc   & $0.0506$ & $0.0916$ & $-1.8660$ & $5.9463$\\
 & cc2  & $0.0504$ & $0.1315$ & $-0.0764$ & $3.1865$\\
 & pp   & $0.0623$ & $0.1593$ & $\phantom{-}1.6231$ & $6.0103$\\
 & pp2  & $0.0621$ & $0.1656$ & $\phantom{-}1.2713$ & $5.2876$\\\midrule
\multirow{4}{*}{4.2 Gaussian $\rho=0.3$}
 & cc   & $0.0507$ & $0.0920$ & $-1.8890$ & $6.0673$\\
 & cc2  & $0.0519$ & $0.2638$ & $\phantom{-}0.1036$ & $3.7720$\\
 & pp   & $0.0622$ & $0.1592$ & $\phantom{-}1.6317$ & $6.0943$\\
 & pp2  & $0.0632$ & $0.2028$ & $\phantom{-}0.5806$ & $3.7156$\\\midrule
\multirow{4}{*}{4.3 $t$-copula $\rho=0.3,\nu=5$}
 & cc   & $0.0511$ & $0.0913$ & $-1.8863$ & $6.0419$\\
 & cc2  & $0.0507$ & $0.2608$ & $-0.0779$ & $4.3228$\\
 & pp   & $0.0623$ & $0.1585$ & $\phantom{-}1.6201$ & $6.0619$\\
 & pp2  & $0.0620$ & $0.2033$ & $\phantom{-}0.7074$ & $4.1131$\\\midrule
\multirow{4}{*}{4.4 Clayton $\theta=10$}
 & cc   & $0.0512$ & $0.0912$ & $-1.8945$ & $6.0986$\\
 & cc2  & $0.0516$ & $0.1425$ & $\phantom{-}0.1637$ & $6.9970$\\
 & pp   & $0.0625$ & $0.1586$ & $\phantom{-}1.6065$ & $5.9304$\\
 & pp2  & $0.0628$ & $0.1596$ & $\phantom{-}1.5584$ & $5.7335$\\
\bottomrule
\end{tabular}
\end{table}

\begin{figure}[H]\centering
\includegraphics[width=0.95\linewidth]{fig_4_4_1_gaussian_high.png}
\caption{4.1 -- Gaussian copula $\rho=0.9$.}\label{fig:4-1}
\end{figure}

\begin{figure}[H]\centering
\includegraphics[width=0.95\linewidth]{fig_4_4_2_gaussian_low.png}
\caption{4.2 -- Gaussian copula $\rho=0.3$.}\label{fig:4-2}
\end{figure}

\begin{figure}[H]\centering
\includegraphics[width=0.95\linewidth]{fig_4_4_3_student_t.png}
\caption{4.3 -- Student-$t$ copula $\rho=0.3$, $\nu=5$.}\label{fig:4-3}
\end{figure}

\begin{figure}[H]\centering
\includegraphics[width=0.95\linewidth]{fig_4_4_4_clayton.png}
\caption{4.4 -- Clayton copula $\theta=10$.}\label{fig:4-4}
\end{figure}

% ---------------------------------------------------------------
\subsection{Discussion}
% ---------------------------------------------------------------

The single-asset strategies \code{cc} and \code{pp} are essentially identical
across the four panels (small Monte-Carlo noise only): they depend on
$S_{1T}$ alone, so the joint distribution of $(Z_1,Z_2)$ does not enter.
The cross-asset strategies \code{cc2} and \code{pp2} react sharply to the
copula choice:

\begin{itemize}\setlength\itemsep{0pt}
\item \emph{Gaussian, high $\rho=0.9$}: \code{cc2} std $=0.13$ -- the held
      asset 2 is almost a copy of asset 1 so the option hedge is almost
      perfect, leaving roughly the same risk profile as \code{cc}.
\item \emph{Gaussian, low $\rho=0.3$}: \code{cc2} std jumps to $0.26$
      because the option payoff (function of $S_{1T}$) decouples from the held
      asset $S_{2T}$.
\item \emph{Student $t$ at the same $\rho=0.3$}: similar std to the Gaussian
      case, but \emph{kurtosis rises from $3.77$ to $4.32$} thanks to the
      tail dependence injected by $\nu=5$. The protective put inherits a
      heavier upper tail (\code{pp2} kurt $4.11$ vs $3.72$).
\item \emph{Clayton $\theta=10$}: lower-tail dependence collapses
      \code{cc2} std back to $0.14$ -- when $S_{1T}$ is small, $S_{2T}$ is
      also small with high probability, so asset 2 plus a long-call
      position behaves like a single-asset covered call. Skew flips
      from $-1.87$ (single asset) to $+0.16$ because the call's
      truncation now no longer captures the asset 2 upside symmetrically.
\end{itemize}
The numbers replicate the slide-pages 16--19 exactly, validating both the
Black--Scholes routine and the copula simulators.
